\documentclass[a4paper]{article}
\usepackage{vntex}
\usepackage{mathptmx}[ptm]
\usepackage{a4wide,amssymb,epsfig,latexsym,array,hhline,fancyhdr}
\usepackage[normalem]{ulem}

\usepackage[makeroom]{cancel}
\usepackage{amsmath}
\usepackage{amsthm}
\usepackage{multicol,longtable,amscd}
\usepackage{diagbox}%Make diagonal lines in tables
\usepackage{booktabs}
\usepackage{alltt}
\usepackage[framemethod=tikz]{mdframed}% For highlighting paragraph backgrounds
\usepackage{caption,subcaption}

\usepackage{lastpage}
\usepackage[lined,boxed,commentsnumbered]{algorithm2e}
\usepackage{enumerate}
\usepackage{color}
\usepackage{graphicx}							% Standard graphics package
\usepackage{array}
\usepackage{tabularx, caption}
\usepackage{multirow}
\usepackage{multicol}
\usepackage{rotating}
\usepackage{graphics}
\usepackage{geometry}
\usepackage{setspace}
\usepackage{epsfig}
\usepackage{tikz}

\usetikzlibrary{arrows,snakes,backgrounds,calc}
\usepackage[unicode]{hyperref}
\hypersetup{urlcolor=blue,linkcolor=black,citecolor=black,colorlinks=true} 
\usepackage{listings}

\usepackage[normalem]{ulem}

\newtheorem{theorem}{{\bf Định lý}}
\newtheorem{property}{{\bf Tính chất}}
\newtheorem{proposition}{{\bf Mệnh đề}}
\newtheorem{corollary}[proposition]{{\bf Hệ quả}}
\newtheorem{lemma}[proposition]{{\bf Bổ đề}}
\theoremstyle{definition}
\newtheorem{exer}{Bài toán}
\addtocontents{toc}{\protect\thispagestyle{empty}}
% Remove page number in contents page. 
\addtocontents{lof}{\protect\thispagestyle{empty}}
% Remove page number in list of figure page. 
\addtocontents{lot}{\protect\thispagestyle{empty}}
% Remove page number in list of table page. 
\def\thesislayout{	% A4: 210 × 297
	\geometry{
		a4paper,
		total={160mm,240mm},  % fix over page
		left=30mm,
		top=22mm,
            right=20mm,
            bottom=20mm,
	}
}
\def\thesisheadlayout{	% A4: 210 × 297
	\geometry{
		a4paper,
		total={160mm,240mm},  % fix over page
		left=30mm,
		top=10mm,
	}
}
\thesislayout
\lstset{
language=R,
basicstyle=\footnotesize\sffamily,
commentstyle=\ttfamily\color{black},
numbers=left,
numberstyle=\ttfamily\color{black}\footnotesize,
stepnumber=1,
numbersep=5pt,
backgroundcolor=\color{white},
showspaces=false,
showstringspaces=false,
showtabs=false,
frame=single,
tabsize=2,
captionpos=b,
breaklines=true,
breakatwhitespace=false,
title=\lstname,
escapeinside={},
keywordstyle={},
morekeywords={}
}
%\usepackage{fancyhdr}
\setlength{\headheight}{40pt}
\pagestyle{fancy}
\fancyhead{} % clear all header fields
\renewcommand{\footruleskip}{1mm}

\fancyhead[L]{
 \begin{tabular}{rl}
    \begin{picture}(25,15)(0,0)
    \put(0,-8){\includegraphics[width=10mm, height=10mm]{images/hcmut.png}}
    %\put(0,-8){\epsfig{width=10mm,figure=hcmut.eps}}
   \end{picture}&
	%\includegraphics[width=8mm, height=8mm]{hcmut.png} & %
	\begin{tabular}{l}
		\textbf{  Trường Đại Học Bách Khoa - Đại học Quốc gia TP.HCM }\\
		\textbf{  Khoa Khoa Học \& Kỹ Thuật Máy Tính}
	\end{tabular} 	
 \end{tabular}
}
\fancyhead[R]{
	\begin{tabular}{l}
		\tiny \bf \\
		\tiny \bf 
	\end{tabular}  }
\fancyfoot{} % clear all footer fields
\fancyfoot[L]{\scriptsize  Báo cáo Bài tập lớn Hướng Công Nghệ Phần Mềm (CO3109) - HK2 2024 - 2025}
\fancyfoot[R]{\scriptsize  Trang {\thepage}/\pageref{LastPage}}

\renewcommand{\headrulewidth}{0.3pt}
\renewcommand{\footrulewidth}{0.3pt}

%%%
\setcounter{secnumdepth}{4}
\setcounter{tocdepth}{3}
\makeatletter
\newcounter {subsubsubsection}[subsubsection]
\renewcommand\thesubsubsubsection{\thesubsubsection .\@alph\c@subsubsubsection}
\newcommand\subsubsubsection{\@startsection{subsubsubsection}{4}{\z@}%
                                     {-3.25ex\@plus -1ex \@minus -.2ex}%
                                     {1.5ex \@plus .2ex}%
                                     {\normalfont\normalsize\bfseries}}
\newcommand*\l@subsubsubsection{\@dottedtocline{3}{10.0em}{4.1em}}
\newcommand*{\subsubsubsectionmark}[1]{}
\makeatother

\everymath{\color{black}}%make in-line maths symbols blue to read/check easily

\sloppy
\captionsetup[figure]{labelfont={small,bf},textfont={small,it},belowskip=-1pt,aboveskip=-9pt}
%space remove between caption, figure, and text
\captionsetup[table]{labelfont={small,bf},textfont={small,it},belowskip=-1pt,aboveskip=7pt}
\setlength{\floatsep}{5pt plus 2pt minus 2pt}
\setlength{\textfloatsep}{5pt plus 2pt minus 2pt}
\setlength{\intextsep}{10pt plus 2pt minus 2pt}

\thesislayout
\onehalfspacing

\begin{document}
\begin{titlepage}
\begin{tikzpicture}[remember picture, overlay]
  \draw[line width = 4pt] ($(current page.north west) + (0.4in,-0.5in)$) rectangle ($(current page.south east) + (-0.4in,0.5in)$);
  \draw[line width=1.5pt]
    ($ (current page.north west) + (0.45in,-0.55in) $)
    rectangle
    ($ (current page.south east) + (-0.45in,0.55in) $);
\end{tikzpicture}

\begin{center}
\LARGE \textbf{ĐẠI HỌC QUỐC GIA THÀNH PHỐ HỒ CHÍ MINH} \\
\vspace{0.2cm}
\LARGE \textbf{TRƯỜNG ĐẠI HỌC BÁCH KHOA} \\
\vspace{0.2cm}
\LARGE \textbf{KHOA KHOA HỌC VÀ KĨ THUẬT MÁY TÍNH}
\end{center}

\vspace{0.3cm}

\begin{figure}[h!]
\begin{center}
\includegraphics[width=4cm]{images/hcmut.png}
\end{center}
\end{figure}

\begin{center}
\begin{tabular}{c}
\textbf{\fontsize{17}{20}\selectfont THỰC TẬP ĐỒ ÁN MÔN HỌC ĐA NGÀNH} \\
\textbf{\fontsize{17}{20}\selectfont HƯỚNG CÔNG NGHỆ PHẦN MỀM (CO3109)} \\[15pt]
\textbf{\fontsize{17}{20}\selectfont ĐỀ TÀI: HỆ THỐNG NHẬN DIỆN KHUÔN MẶT} \\
\textbf{\fontsize{17}{20}\selectfont CHO ĐIỂM DANH VÀ QUẢN LÝ CỬA HÀNG} \\[15pt]
\\
\end{tabular}
\end{center}

\vspace{0.5cm}
\begin{table}[h]
\begin{tabular}{rll}
\hspace{2 cm} &   \textbf{\Large GVHD:} {\Large Nguyễn Thành Lộc}
\\
\\
\\
\hspace{3 cm} &     \textbf{$\;$ $\;$$\;$$\;$$\;$$\;$$\;$$\;$$\;$$\;$$\;$$\;$$\;$$\;$$\;$$\;$$\;$$\;$$\;$$\;$$\;$$\;$$\;$$\;$$\;$$\;$$\;$$\;$$\;$$\;$$\;$$\;$$\;$ \Large ---o0o---} 
\\
\\
\hspace{3 cm} &   	\textbf{\Large SVTH:  } {\Large Nguyễn Tấn Tài - 2212990}
\\
\\
\hspace{3 cm} &   	\hspace{1.65cm} {\Large Nguyễn Văn Tuệ - 2212345}
\\
\\
\end{tabular}
\end{table}
\vspace{0.5cm}
\begin{center}
{\Large TP. HỒ CHÍ MINH, THÁNG 4/2025 }
\end{center}
\end{titlepage}
%\thispagestyle{empty}
\newpage
%%%%%%%%%%%%%%%%%INTRO%%%%%%%%%%%%%%%%%%%
\newpage
\tableofcontents
\newpage
\listoffigures
\newpage
\listoftables
\newpage
\setcounter{page}{1}

%%%%%%%%%%%%%%%%%CONTENTS%%%%%%%%%%%%%%%%%%%
\section{Giới thiệu}

\subsection{Bối cảnh nghiên cứu}
\hspace{0.5cm}Ô nhiễm không khí là một trong những vấn đề môi trường nghiêm trọng nhất hiện nay, đặc biệt là tại các khu vực đô thị. Trong đó, Carbon Monoxide (CO) và Carbon Dioxide (CO2) là hai trong số các chất ô nhiễm không khí quan trọng cần được theo dõi và kiểm soát. CO là một khí độc không màu, không mùi, có thể gây nguy hiểm đến sức khỏe con người khi tiếp xúc ở nồng độ cao. CO2, mặc dù ít độc hơn, nhưng là một trong những khí nhà kính chính góp phần vào biến đổi khí hậu toàn cầu.

\subsection{Mục tiêu nghiên cứu}
\hspace{0.5cm}Dự án này tập trung vào việc phát triển các mô hình học máy để dự đoán nồng độ CO và CO2 trong không khí dựa trên bộ dữ liệu Air Quality từ UCI Machine Learning Repository. Cụ thể, các mục tiêu chính bao gồm:

\begin{itemize}
    \item Phân tích và xử lý dữ liệu từ các cảm biến chất lượng không khí
    \item Xây dựng các mô hình dự đoán nồng độ CO và CO2
    \item Đánh giá hiệu suất của các mô hình dự đoán
    \item Phân tích tầm quan trọng của các yếu tố ảnh hưởng đến nồng độ CO và CO2
\end{itemize}

\subsection{Bộ dữ liệu}
\hspace{0.5cm}Bộ dữ liệu Air Quality được thu thập từ một thiết bị đa cảm biến khí được triển khai tại một thành phố ở Ý. Dữ liệu bao gồm:

\begin{itemize}
    \item Thời gian thu thập: Từ tháng 3/2004 đến tháng 2/2005
    \item Số lượng mẫu: 9,358 mẫu
    \item Tần suất: Trung bình theo giờ
    \item Các đặc trưng chính:
    \begin{itemize}
        \item Nồng độ CO (mg/m³)
        \item Nồng độ các hợp chất hữu cơ không chứa metan (µg/m³)
        \item Nồng độ Benzen (µg/m³)
        \item Nồng độ NOx (ppb)
        \item Nồng độ NO2 (µg/m³)
        \item Nhiệt độ (°C)
        \item Độ ẩm tương đối (\%)
        \item Độ ẩm tuyệt đối (g/m³)
    \end{itemize}
\end{itemize}

\subsection{Ý nghĩa thực tiễn}
\hspace{0.5cm}Nghiên cứu này có ý nghĩa quan trọng trong việc:

\begin{itemize}
    \item Cảnh báo sớm về mức độ ô nhiễm không khí
    \item Hỗ trợ các cơ quan quản lý môi trường trong việc ra quyết định
    \item Nâng cao nhận thức cộng đồng về chất lượng không khí
    \item Góp phần vào việc phát triển các hệ thống giám sát chất lượng không khí thông minh
\end{itemize}

\subsection{Cấu trúc báo cáo}
\hspace{0.5cm}Báo cáo được tổ chức thành các phần chính sau:

\begin{itemize}
    \item Phần 1: Giới thiệu tổng quan về đề tài
    \item Phần 2: Cơ sở lý thuyết về học máy và các mô hình dự đoán
    \item Phần 3: Phương pháp nghiên cứu và xử lý dữ liệu
    \item Phần 4: Kết quả thực nghiệm và đánh giá
    \item Phần 5: Kết luận và hướng phát triển
\end{itemize}


\newpage
%%%%%%%%%%%%%%%%%OUTRO%%%%%%%%%%%%%%%%%%%
\section{Tài liệu tham khảo}

\begin{thebibliography}{9}

\bibitem{phamkhanh}
Phạm Đình Khánh.
\textit{Deep AI Book - Random Forest}.
\url{https://phamdinhkhanh.github.io/deepai-book/ch_ml/RandomForest.html}

\bibitem{tabml}
Machine Learning Cơ Bản.
\textit{Random Forest}.
\url{https://machinelearningcoban.com/tabml_book/ch_model/random_forest.html}

\bibitem{geeksforgeeks_xgboost}
GeeksforGeeks.
\textit{XGBoost}.
\url{https://www.geeksforgeeks.org/xgboost/}

\bibitem{viblo_ann}
Viblo.
\textit{Tổng quan về Artificial Neural Network}.
\url{https://viblo.asia/p/tong-quan-ve-artificial-neural-network-1VgZvwYrlAw}

\bibitem{viblo_lr}
Viblo.
\textit{Linear Regression - Hồi quy tuyến tính trong Machine Learning}.
\url{https://viblo.asia/p/linear-regression-hoi-quy-tuyen-tinh-trong-machine-learning-4P856akRlY3}

\bibitem{geeksforgeeks_lr}
GeeksforGeeks.
\textit{ML Linear Regression}.
\url{https://www.geeksforgeeks.org/ml-linear-regression/}

\end{thebibliography}
\end{document}
